\section{Local Rules A/R/G/M and Static Guard Compilation}
\label{app:guard-rules}

Section~\ref{sec:guard-exactification} states the structural compilation theorem used by the exactification argument. This appendix supplies the local rule-level proofs. Its role is deliberately narrow. We do \emph{not} yet prove aligned-leaf exactness, global split-tree cover assembly, or finite exactification of the whole network; those belong to Appendix~\ref{app:exactification}. The present appendix proves only that the static compilation manipulates finite verifier-visible objects and that each local rule preserves certified faithfulness.

The basic formal backend for certified cells and payloads was fixed in Appendix~\ref{app:cell-payload-syntax}. Here we specialize it to the exact side of the guard language. Throughout the appendix:
\begin{itemize}
    \item $D\subseteq \mathbb R^n$ is nonempty compact;
    \item $u$ denotes a scalar node in an allowed CPWL expression graph;
    \item $[[u]]:D\to\mathbb R$ denotes the true scalar semantics of that node;
    \item the guard language is generated by affine comparison literals.
\end{itemize}

Allowed CPWL expression graphs are assumed to be already normalized so that the affine primitives are exactly unary affine post-composition and addition. Any finite-input affine map can be decomposed into this basis by introducing finitely many scalar multiples and additions, so the compilation rules below lose no generality.

For every affine function
\[
 h(x)=w^\top x+b,
\]
we write the associated closed-halfspace literals as
\[
 h\le 0
 \qquad\text{and}\qquad
 h\ge 0.
\]
A finite set of such literals is a \emph{guard set} $S$, with semantic cell
\[
 D_S:=D\cap \bigcap_{\lambda\in S}\lambda.
\]
When $S$ and $T$ are guard sets, their common refinement cell is simply
\[
 D_{S\cup T}=D_S\cap D_T.
\]

A \emph{certified guard-affine fragment} for a scalar node $u$ is a triple
\[
 (S,a,\Pi),
\]
where $S$ is a guard set, $a(x)=w^\top x+b$ is affine, and $\Pi$ is a finite payload accepted by the verifier as a proof of
\[
 [[u]](x)=a(x)
 \qquad \forall x\in D_S.
\]
A finite family
\[
 \mathcal F_u=\{(S_r,a_r,\Pi_r)\}_{r=1}^{R_u}
\]
is \emph{certified faithful} for $u$ if:
\begin{enumerate}
    \item the cells $\{D_{S_r}\}_{r=1}^{R_u}$ cover $D$; and
    \item every payload $\Pi_r$ is accepted as a proof of
    \[
    [[u]](x)=a_r(x)
    \qquad \forall x\in D_{S_r}.
    \]
\end{enumerate}
Thus every compiled fragment is already an accepted local exact statement in the sense of Section~\ref{sec:verifier-model}; what later appendices add is leaf alignment and a global cover payload.

\subsection{Preliminaries: Restriction and Common Refinement}

The rule proofs repeatedly use two elementary closure principles.

\paragraph{Restriction of accepted exact payloads.}
Suppose the verifier accepts a finite payload proving
\[
 [[u]](x)=a(x)
 \qquad \forall x\in D_S,
\]
and also accepts a finite containment proof
\[
 D_{S'}\subseteq D_S.
\]
Then the same equality is valid on the smaller cell $D_{S'}$. We therefore obtain a derived finite payload, denoted
\[
 \mathrm{Restr}(S'\subseteq S;\Pi),
\]
that the verifier accepts as a proof of
\[
 [[u]](x)=a(x)
 \qquad \forall x\in D_{S'}.
\]
This is exactly the downward-closure rule for accepted exact payloads already implicit in Appendix~\ref{app:cell-payload-syntax}.

\paragraph{Common refinement of certified faithful families.}
If
\[
 \mathcal F_1=\{(S_i,a_i,\Pi_i)\}_{i=1}^{R_1}
 \qquad\text{and}\qquad
 \mathcal F_2=\{(T_j,b_j,\Theta_j)\}_{j=1}^{R_2}
\]
both cover $D$, then the refined family of cells
\[
 \{D_{S_i\cup T_j}\}_{1\le i\le R_1,\ 1\le j\le R_2}
\]
also covers $D$. Indeed, for every $x\in D$ there exist indices $i,j$ with $x\in D_{S_i}$ and $x\in D_{T_j}$, hence
\[
 x\in D_{S_i\cup T_j}=D_{S_i}\cap D_{T_j}.
\]
We use this observation in Rules~R and~M.

\subsection{Rule A: Affine Propagation}

Assume a scalar node $y$ is obtained from a scalar node $\phi$ by affine post-composition:
\[
 y(x)=\alpha\phi(x)+\beta.
\]
Suppose
\[
 \mathcal F_\phi=\{(S_r,a_r,\Pi_r)\}_{r=1}^{R}
\]
is certified faithful for $\phi$. Rule~A outputs the family
\[
 \mathcal F_y:=\{(S_r,\alpha a_r+\beta,\Pi_r^{\mathrm A})\}_{r=1}^{R},
\]
where each new payload $\Pi_r^{\mathrm A}$ is the finite verifier-visible record consisting of:
\begin{enumerate}
    \item the already accepted child payload $\Pi_r$;
    \item the affine coefficients $\alpha,\beta$;
    \item the declaration that the new node is obtained from the child node by affine propagation.
\end{enumerate}
The checker only needs to confirm that the guard set has not changed and that the new affine tag is exactly the symbolic transformation of the old one.

\paragraph{Theorem C.1 (Rule A).}
Rule~A preserves certified faithfulness.

\paragraph{Proof.}
Coverage is immediate, since the output family uses exactly the same guard cells $D_{S_r}$ as the input family.

Fix one input fragment $(S_r,a_r,\Pi_r)$ and one point $x\in D_{S_r}$. Because $\Pi_r$ is accepted, the verifier has already established
\[
 [[\phi]](x)=a_r(x).
\]
Therefore
\[
 [[y]](x)=\alpha[[\phi]](x)+\beta=\alpha a_r(x)+\beta.
\]
So the proposed output tag is semantically correct on $D_{S_r}$.

It remains to justify verifier acceptance. The payload $\Pi_r^{\mathrm A}$ is finite and introduces no new semantic search object: it only points to an already accepted equality payload, records the affine post-composition coefficients, and asks the checker to validate a one-step symbolic transformation. Hence the verifier accepts
\[
 [[y]](x)=\alpha a_r(x)+\beta
 \qquad \forall x\in D_{S_r}.
\]
Thus every output fragment is locally exact, and the family remains certified faithful. \qed

\subsection{Rule R: Addition on a Common Refinement}

Assume
\[
 y(x)=\phi_1(x)+\phi_2(x),
\]
and suppose
\[
 \mathcal F_1=\{(S_i,a_i,\Pi_i)\}_{i=1}^{R_1},
 \qquad
 \mathcal F_2=\{(T_j,b_j,\Theta_j)\}_{j=1}^{R_2}
\]
are certified faithful for $\phi_1$ and $\phi_2$, respectively. Rule~R outputs the common-refinement family
\[
 \mathcal F_y:=\bigl\{(S_i\cup T_j,\ a_i+b_j,\ \Pi_{ij}^{\mathrm R}) : 1\le i\le R_1,\ 1\le j\le R_2\bigr\},
\]
where $\Pi_{ij}^{\mathrm R}$ is the finite payload recording:
\begin{enumerate}
    \item the already accepted child payloads $\Pi_i$ and $\Theta_j$;
    \item the new guard set $S_i\cup T_j$;
    \item the symbolic identity that the new affine tag is the sum $a_i+b_j$.
\end{enumerate}

\paragraph{Theorem C.2 (Rule R).}
Rule~R preserves certified faithfulness.

\paragraph{Proof.}
Coverage follows from common refinement: because the two input families both cover $D$, the refined cells $D_{S_i\cup T_j}$ cover $D$ as $(i,j)$ ranges over all pairs.

Now fix an output fragment
\[
 (S_i\cup T_j,\ a_i+b_j,\ \Pi_{ij}^{\mathrm R})
\]
and a point $x\in D_{S_i\cup T_j}$. Since
\[
 D_{S_i\cup T_j}\subseteq D_{S_i}\cap D_{T_j},
\]
the already accepted child payloads imply
\[
 [[\phi_1]](x)=a_i(x),
 \qquad
 [[\phi_2]](x)=b_j(x).
\]
Therefore
\[
 [[y]](x)= [[\phi_1]](x)+[[\phi_2]](x)=a_i(x)+b_j(x).
\]
So the proposed affine tag is semantically exact on the refined cell.

Verifier acceptance is again finite and local. The checker only has to confirm that the child payloads are already accepted, that the new cell is the common-refinement cell $D_{S_i\cup T_j}$, and that the announced affine tag is the symbolic sum of the two child tags. Hence it accepts the local exact statement on $D_{S_i\cup T_j}$. The output family is therefore certified faithful. \qed

\subsection{Rule G: Gating by an Affine Winner Comparison}

Let $y=\sigma(\phi)$, where $\sigma$ is one of the scalar gates used in Section~\ref{sec:guard-exactification}:
\[
 \operatorname{ReLU}(z)=\max\{z,0\},
 \qquad
 \operatorname{Leaky}_\alpha(z)=\max\{z,\alpha z\},
 \qquad
 \operatorname{Abs}(z)=\max\{z,-z\}.
\]
We also allow PReLU with a fixed compile-time parameter $\alpha$, handled exactly like the Leaky case. The standard Leaky/PReLU regime is $0\le \alpha\le 1$, but the rule below is correct for every fixed compile-time $\alpha\in\mathbb R$.

For a fixed gate type $\sigma$, define the two affine branch templates
\[
 b_1^\sigma(z):=z,
 \qquad
 b_2^\sigma(z):=
 \begin{cases}
  0, & \text{for ReLU},\\
  \alpha z, & \text{for Leaky/PReLU},\\
  -z, & \text{for Abs},
 \end{cases}
\]
and their affine winner comparison
\[
 d_\sigma(z):=b_1^\sigma(z)-b_2^\sigma(z).
\]
Thus
\[
 \sigma(z)=\max\{b_1^\sigma(z),b_2^\sigma(z)\},
\]
and the active affine branch is decided by the sign of $d_\sigma(z)$.

For an input fragment $(S,a,\Pi)$, write
\[
 a(x)=w^\top x+b,
\]
and abbreviate the corresponding affine branch tags by
\[
 b_1(a):=b_1^\sigma\!\circ a,
 \qquad
 b_2(a):=b_2^\sigma\!\circ a,
 \qquad
 d(a):=d_\sigma\!\circ a=b_1(a)-b_2(a).
\]
Rule~G now branches on the actual winner comparison $d(a)$:
\begin{enumerate}
    \item If $d(a)$ is nonconstant, the rule outputs the two fragments
    \[
    (S\cup\{d(a)\ge 0\},\ b_1(a),\ \Pi^{\mathrm G,1}),
    \qquad
    (S\cup\{d(a)\le 0\},\ b_2(a),\ \Pi^{\mathrm G,2}).
    \]
    \item If $d(a)$ is a nonnegative constant, the first branch wins everywhere on $D_S$, so the rule outputs the single fragment
    \[
    (S,\ b_1(a),\ \Pi^{\mathrm G,\mathrm{const}+}).
    \]
    \item If $d(a)$ is a negative constant, the second branch wins everywhere on $D_S$, so the rule outputs the single fragment
    \[
    (S,\ b_2(a),\ \Pi^{\mathrm G,\mathrm{const}-}).
    \]
\end{enumerate}
In the degenerate case $d(a)\equiv 0$, the two affine branches coincide identically, so either constant-branch output is acceptable.

All Rule~G payloads are finite. Each records:
\begin{enumerate}
    \item the accepted child payload $\Pi$;
    \item the optional new winner literal when $d(a)$ is nonconstant;
    \item a branch tag recording which affine branch $b_1(a)$ or $b_2(a)$ is claimed to win on the current cell; and
    \item the finite symbolic formula for the chosen affine branch and its comparison $d(a)$.
\end{enumerate}
The checker only has to verify that the announced branch formulas match the gate type and that the literal being added is exactly the affine comparison deciding between those two branches.

\paragraph{Theorem C.3 (Rule G).}
Rule~G preserves certified faithfulness.

\paragraph{Proof.}
Fix one input fragment $(S,a,\Pi)$. Because $\Pi$ is accepted,
\[
 [[\phi]](x)=a(x)
 \qquad \forall x\in D_S.
\tag{C.1}
\]
We prove coverage and local exactness case by case.

\emph{Nonconstant-comparison case.}
Assume $d(a)$ is nonconstant. Since the two closed halfspaces $\{d(a)\le 0\}$ and $\{d(a)\ge 0\}$ cover all of $\mathbb R^n$, we have
\[
 D_S\subseteq D_{S\cup\{d(a)\le 0\}}\ \cup\ D_{S\cup\{d(a)\ge 0\}}.
\]
Thus replacing the input fragment by the two output fragments preserves coverage.

Take any point $x\in D_{S\cup\{d(a)\ge 0\}}$. Then $x\in D_S$ and $d(a)(x)\ge 0$, so
\[
 b_1(a)(x)\ge b_2(a)(x).
\]
Using (C.1),
\[
 [[y]](x)=\sigma([[\phi]](x))=\sigma(a(x))=\max\{b_1(a)(x),b_2(a)(x)\}=b_1(a)(x).
\]
Hence the first output fragment is semantically exact on its cell.

Similarly, for any point $x\in D_{S\cup\{d(a)\le 0\}}$, the inequality $d(a)(x)\le 0$ implies
\[
 b_2(a)(x)\ge b_1(a)(x),
\]
so again by (C.1),
\[
 [[y]](x)=\sigma(a(x))=\max\{b_1(a)(x),b_2(a)(x)\}=b_2(a)(x).
\]
Thus the second output fragment is semantically exact on its cell.

The overlap of the two child cells is contained in
\[
 \{x:d(a)(x)=0\}.
\]
On that overlap the two branch tags agree because
\[
 b_1(a)(x)-b_2(a)(x)=d(a)(x)=0.
\]
Therefore boundary points cause no consistency problem: if both branch literals hold, both local affine equalities certify the same semantic value.

\emph{Constant-comparison case.}
Assume now that $d(a)$ is constant. If $d(a)\ge 0$, then $b_1(a)(x)\ge b_2(a)(x)$ for every $x\in D_S$, so
\[
 [[y]](x)=\sigma(a(x))=b_1(a)(x)
 \qquad \forall x\in D_S.
\]
Hence the single output fragment $(S,b_1(a),\Pi^{\mathrm G,\mathrm{const}+})$ is locally exact and coverage is unchanged. If instead $d(a)<0$, the same argument gives
\[
 [[y]](x)=b_2(a)(x)
 \qquad \forall x\in D_S,
\]
so the single output fragment $(S,b_2(a),\Pi^{\mathrm G,\mathrm{const}-})$ is locally exact. When $d(a)\equiv 0$, the two affine branches coincide identically and either single-fragment output is valid.

In all cases every output fragment is locally exact, all payloads remain finite, and the union of the output cells covers the union of the input cells. Thus Rule~G preserves certified faithfulness. \qed

\subsection{Rule M: Finite Max with Explicit Winner Constraints}

Assume
\[
 y(x)=\max_{1\le i\le k} \phi_i(x),
\]
and for each child $\phi_i$ suppose we have a certified faithful family
\[
 \mathcal F_i=\{(S_{i,r},a_{i,r},\Pi_{i,r})\}_{r=1}^{R_i}.
\]
For every index tuple
\[
 \mathbf r=(r_1,\dots,r_k),
\]
define the common-refinement guard set
\[
 S_{\mathbf r}:=\bigcup_{i=1}^k S_{i,r_i}
\]
and abbreviate the corresponding child affine tags by
\[
 a_i:=a_{i,r_i}
 \qquad (1\le i\le k).
\]
Thus on the cell $D_{S_{\mathbf r}}$, every child node is already represented by an accepted affine tag.

Now fix a candidate winner index $p\in\{1,\dots,k\}$. For every competitor $q\neq p$, define the affine comparison
\[
 d_{pq}(x):=a_p(x)-a_q(x).
\]
Rule~M treats each competitor as follows.
\begin{enumerate}
    \item If $d_{pq}$ is a negative constant, then $p$ can never dominate $q$ on this common refinement, so the candidate $p$ is discarded.
    \item If $d_{pq}$ is a nonnegative constant, then no new literal is needed; the comparison is globally valid on the refinement cell.
    \item If $d_{pq}$ is nonconstant, the rule adds the winner literal
    \[
    d_{pq}\ge 0.
    \]
\end{enumerate}
Let $S(\mathcal C_{p,\mathbf r})$ denote the finite set of all nonconstant winner literals added for the surviving candidate $p$ on the tuple $\mathbf r$. If $p$ is not discarded, Rule~M outputs the fragment
\[
 \bigl(S_{\mathbf r}\cup S(\mathcal C_{p,\mathbf r}),\ a_p,\ \Pi_{p,\mathbf r}^{\mathrm M}\bigr).
\]
The payload $\Pi_{p,\mathbf r}^{\mathrm M}$ is finite and contains:
\begin{enumerate}
    \item the tuple of already accepted child payloads $\Pi_{1,r_1},\dots,\Pi_{k,r_k}$;
    \item the chosen winner index $p$;
    \item for each competitor $q\neq p$, either the explicit constant comparison $d_{pq}\ge 0$ or the newly added literal $d_{pq}\ge 0$;
    \item the declaration that the output affine tag is the winner tag $a_p$.
\end{enumerate}

\paragraph{Theorem C.4 (Rule M).}
Rule~M preserves certified faithfulness.

\paragraph{Proof.}
We must prove both coverage and local exactness.

\emph{Coverage.}
Fix any point $x\in D$. Because each child family $\mathcal F_i$ is certified faithful, for every $i$ there exists some $r_i$ such that
\[
 x\in D_{S_{i,r_i}}.
\]
Hence for the tuple $\mathbf r=(r_1,\dots,r_k)$ we have
\[
 x\in D_{S_{\mathbf r}}.
\]
On this refinement cell, each child semantics is represented by the corresponding affine tag:
\[
 [[\phi_i]](x)=a_i(x)
 \qquad (1\le i\le k).
\]
Now choose any winner index
\[
 p^\star\in \operatorname*{arg\,max}_{1\le i\le k} a_i(x).
\]
For every competitor $q\neq p^\star$, the comparison
\[
 a_{p^\star}(x)-a_q(x)\ge 0
\]
holds. Therefore:
\begin{itemize}
    \item if $d_{p^\star q}$ is a constant, it cannot be negative, otherwise $p^\star$ would fail to dominate $q$ at $x$;
    \item if $d_{p^\star q}$ is nonconstant, then $x$ satisfies the corresponding winner literal $d_{p^\star q}\ge 0$.
\end{itemize}
Thus $x$ belongs to the surviving winner fragment for $p^\star$:
\[
 x\in D_{S_{\mathbf r}\cup S(\mathcal C_{p^\star,\mathbf r})}.
\]
So the Rule~M output family still covers $D$.

\emph{Local exactness.}
Fix one surviving output fragment
\[
 \bigl(S_{\mathbf r}\cup S(\mathcal C_{p,\mathbf r}),\ a_p,\ \Pi_{p,\mathbf r}^{\mathrm M}\bigr)
\]
and a point
\[
 x\in D_{S_{\mathbf r}\cup S(\mathcal C_{p,\mathbf r})}.
\]
Since $x\in D_{S_{\mathbf r}}$, the child payloads give
\[
 [[\phi_i]](x)=a_i(x)
 \qquad (1\le i\le k).
\]
Moreover, for every competitor $q\neq p$, either:
\begin{itemize}
    \item $d_{pq}$ is a nonnegative constant, hence $a_p(x)\ge a_q(x)$ automatically; or
    \item the literal $d_{pq}\ge 0$ is part of the guard set, hence again $a_p(x)\ge a_q(x)$ on the current cell.
\end{itemize}
Therefore $a_p(x)$ dominates all competitors and
\[
 [[y]](x)=\max_{1\le i\le k}[[\phi_i]](x)=\max_{1\le i\le k} a_i(x)=a_p(x).
\]
So the proposed winner tag is semantically exact on its output cell.

The remaining issue is overlap consistency. Different winner fragments may overlap only where several affine competitors tie. But on every such overlap, all active winner tags take the same value, since they are equal at the overlap point by the relevant winner inequalities in both directions. Thus local exactness is compatible across overlapping winner cells.

Finally, verifier acceptance is finite. The payload $\Pi_{p,\mathbf r}^{\mathrm M}$ only points to already accepted child payloads, records a finite tuple of affine comparisons, and asks the checker to verify that each required winner condition is either a stored nonnegative constant fact or a literal explicitly added to the guard set. Hence the checker accepts the announced exact equality on the output cell. Therefore Rule~M preserves certified faithfulness. \qed

\subsection{Topological Induction over the Expression Graph}

We now combine the four local closure rules into the global static compilation theorem for allowed CPWL expression graphs.

\paragraph{Base fragments.}
Input coordinates and scalar constants give immediate one-fragment certified faithful families:
\[
 \mathcal F_{x_p}=\{(\varnothing,\ x\mapsto x_p,\ \Pi_{x_p}^{\mathrm{ax}})\},
 \qquad
 \mathcal F_c=\{(\varnothing,\ x\mapsto c,\ \Pi_c^{\mathrm{ax}})\}.
\]
Here $\Pi_{x_p}^{\mathrm{ax}}$ and $\Pi_c^{\mathrm{ax}}$ are the trivial finite atomic payloads accepted by the verifier for the corresponding identities on all of $D$.

\paragraph{Theorem C.5 (topological-induction static compilation theorem).}
For every scalar node $u$ in an allowed CPWL expression graph, there exists a finite certified faithful family
\[
 \mathcal F_u=\{(S_r,a_r,\Pi_r)\}_{r=1}^{R_u}
\]
of guard-affine fragments such that each $\Pi_r$ is accepted as a proof of
\[
 [[u]](x)=a_r(x)
 \qquad \forall x\in D_{S_r},
\]
and the cells $\{D_{S_r}\}_{r=1}^{R_u}$ cover $D$.

\paragraph{Proof.}
Order the nodes of the finite expression graph topologically.

\emph{Base case.}
For input coordinates and constants, the base fragment families above are already certified faithful.

\emph{Inductive step.}
Assume certified faithful families have been constructed for all predecessors of the current node $u$. There are only four possible module types:
\begin{enumerate}
    \item if $u$ is obtained by affine post-composition from one predecessor, apply Rule~A and invoke Theorem~C.1;
    \item if $u$ is an addition node, apply Rule~R and invoke Theorem~C.2;
    \item if $u$ is a scalar gate, apply Rule~G and invoke Theorem~C.3;
    \item if $u$ is a finite-max node, apply Rule~M and invoke Theorem~C.4.
\end{enumerate}
In each case, the resulting output family is again certified faithful. Therefore the induction proceeds through the whole DAG and yields a certified faithful family for every scalar node, in particular for the output node $u_g$. \qed

\paragraph{Corollary C.6 (finiteness of compilation hyperplanes).}
For every scalar node $u$, the certified faithful family $\mathcal F_u$ introduces only finitely many affine comparison hyperplanes. In particular, the output-node family $\mathcal F_g$ determines a finite comparison-hyperplane set
\[
 \mathcal H^\star.
\]

\paragraph{Proof.}
The base fragments introduce none. Rule~A introduces no new hyperplanes. Rule~R only takes common refinements of already finite guard sets. Rule~G introduces at most one new gate-comparison hyperplane per input fragment, and Rule~M introduces only finitely many pairwise winner-comparison hyperplanes because each max node has finitely many children and each child family is finite. Induction over the finite expression graph therefore yields only finitely many hyperplanes overall. \qed

This completes the purely local part of the CPWL theorem package. We now know that static compilation produces only finitely many verifier-visible guard-affine fragments, each already carrying an accepted exact local payload. Appendix~\ref{app:exactification} uses the finite hyperplane set $\mathcal H^\star$ to build aligned split trees, restricts these compiled fragments to nonempty leaves, and assembles the resulting exact leaves into a finite accepted exact certified cover.
